\section{Foundations of Secure E2E-V Voting Protocols}
The goal of end-to-end verifiable (E2E-V) voting protocols is to guarantee the integrity of an election that can be \emph{verified} by voters and observers alike, all without relying on a central authority \cite{Cha02,kremer2010election,Ben87}.
E2E verifiability requires two levels of verification: 

\begin{itemize}
    \item \textbf{Individual Verifiability:} Each voter can verify their own vote was correctly included in the election.
    \item \textbf{Universal Verifiability:} Any observer can verify the overall election outcome is correct based on the recorded votes \cite{DBLP:journals/corr/AliM16}.
\end{itemize}

To achieve E2E verifiability, we require three core properties: \emph{cast-as-intended}, \emph{recorded-as-cast}, and \emph{counted-as-collected} \cite{regenscheid-e2e-vvsg2-tgdc-2023}.
Together, these properties give confidence to voters that their vote was captured and counted correctly in the election.  
If any error or tampering occurs at any stage of the election, it will be detected by the voter or auditors. 

\subsubsection{Cast-as-Intended}
A voting protocol satisfies this property if every voter can independently confirm whoever they \emph{intended} to vote for is accurately reflected on their ballot \cite{10.1007/978-3-662-63958-0_22}. 
For electronic voting protocols, this translates to their ballot being correctly constructed by the software or voting machine.
This is especially important in scenarios where the voting machine may be compromised or malicious.
In such cases, the machine might display \emph{Candidate\,A} on the screen, but actually encrypt a vote for \emph{Candidate\,B}.
We describe this property formally as follows:

\begin{definition}
Let $\mathcal{V}$ be the voter's intended vote and $\mathcal{B}$ be the voter's ballot. 
A voting protocol provides cast-as-intended verifiability iff there exists a verification process $\text{Verify}_{\text{CAI}}$ such that:
\begin{equation*}
\Pr[\mathsf{Verify}_{CAI}(\mathcal{B}, \mathcal{V}) = 1 \mid \mathcal{B} \text{ encodes } \mathcal{V}] \geq 1 - \text{negl}(\lambda)
\end{equation*}
\begin{equation*}
\Pr[\mathsf{Verify}_{CAI}(\mathcal{B}, \mathcal{V}) = 1 \mid \mathcal{B} \text{ does not encode } \mathcal{V}] \leq \text{negl}(\lambda)
\end{equation*}
where $\lambda$ is the security parameter and $\text{negl}(\lambda)$ denotes a negligible function.
\end{definition}

\subsubsection{Recorded-as-Cast}
A voting protocol satisfies this property if every voter can independently confirm that their cast ballot has been correctly \emph{recorded} in the public election record.
In typical E2E-V voting protocols, this is achieved by providing the voter with a receipt that contains a \emph{tracking code} for their ballot that can be used to check a public list of ballot records known as the \emph{public bulletin board} \cite{guasch2016challenge}.
The bulletin board must be \textit{publicly verifiable} 
This ensures once a ballot is recorded, it cannot be deleted or modified without detection \cite{CRST15}.
If the code appears in the bulletin board, the voter is assured their vote is included in the election.
We describe this property formally as follows:

\begin{definition}
Let $\mathcal{BB}$ denote the public bulletin board and $r$ be the voter's receipt. 
A voting protocol provides recorded-as-cast verifiability iff:
\begin{equation*}
r \in \mathcal{BB} \rightarrow \exists \mathcal{B} \in \mathcal{BB}\,\, s.t. \,\,H(\mathcal{B}) = r
\end{equation*}
where $H$ is a cryptographic hash function.
\end{definition}

\subsubsection{Counted-as-Collected}
A voting protocol satisfies this property if any observer can verify that all collected ballots are \emph{correctly included} in the final tally. 
In typical E2E-V protocols, this is achieved by using cryptographic tools such as homomorphic encryption (where encrypted values can be summed without decryption), or mix-nets (where ballots are shuffled and re-encrypted).
The protocol then provides some proofs to show the correctness of the tally, all without revealing individual votes \cite{HS00}.
We describe this property formally as follows:

\begin{definition}
Let $\mathcal{T}$ denote the published tally, $\mathcal{BB} = \{\mathcal{B}_1, \ldots, \mathcal{B}_n\}$ be the set of ballots on the bulletin board, and $\pi$ be a cryptographic proof of correct tallying. 
A voting protocol provides counted-as-collected verifiability iff there exists a publicly computable function $\mathsf{Verify}_{CAC}$ such that:
\begin{equation*}
\mathsf{Verify}_{CAC}(\mathcal{T}, \mathcal{BB}, \pi) = true \,\, s.t. \,\,\mathcal{T} = \text{Tally}(\mathcal{BB})
\end{equation*}
\end{definition}

\subsection{Levels of Ballot Secrecy}

Along with verifiability, E2E-V voting protocols must also satisfy strong privacy and security properties to protect voters against vote-selling and coercion.
These core properties are \emph{ballot secrecy}, \emph{receipt freeness}, and \emph{coercion resistance}.
Each of these notions are important for free and fair elections. 

\subsubsection{Ballot Secrecy}
Ballot secrecy (also known as voter anonymity or privacy) requires that no party learns how an individual voter voted, beyond what can be inferred from the overall outcome.
In this notion, the adversary is assumed to be \emph{passive}. 
It should not be able to link a cast ballot back to it's voter from the from the observable.
Josh Benaloh, a senior cryptographer at Microsoft Research emphasized "A secret ballot requires that \dots no voter's choices can be surreptitiously released or inferred. The goal of a secret ballot is to provide an election free from undue influence" \cite{Ben06}.
Ballot secrecy is implemented by encrypting or anonymizing ballots, and ensuring any published results do not reveal individual vote contents.
However, ballot secrecy alone does not prevent coercion or vote-selling, as a voter might still be able to prove their vote to a third party. 

\subsubsection{Receipt Freeness}
Receipt freeness is a stronger notion that requires a voter to not be able to obtain any proof or \emph{receipt} that would credibly show how they voted to a third party.
Here, the assumption is that the adversary is post-factum, where it can \emph{check in} with the voter after the election and ask for evidence of how they voted. 
This property was first introduced by Benaloh and Tuinstra to prevent vote-selling and coercion \cite{BT94}.
Ron Rivest, a pioneer in cryptographic voting, emphasized "the voter must not be given a 'receipt' that would allow them to prove how they voted to someone else-otherwise the voter could be coerced or bribed" \cite{RS07}.
Receipt freeness is implemented by ensuring that any data a voter can retain does not allow them to prove how they voted.
Additionally, it ensures no voter has access to the secret \emph{randomness} used to encrypt their vote. 
However, receipt freeness does not protect against more powerful coercion attacks where the adversary can pressure the voter during the voting process itself.

\subsubsection{Coercion Resistance}
Coercion resistance is the strongest of the three properties, requiring even if a voter is coerced or bribed during the voting process, they can still vote freely without the coercer being able to detect it.
In this notion, the adversary is \emph{interactive} and \emph{omnipresent}, where it can stand over the voter's shoulder, provide the encryption keys itself, or force the voter to share their screen.
This property was introduced by Juels, Catalano, and Jakobsson to model strong coercion threats.
They stated "We define a scheme to be coercion-resistant if it is infeasible for the adversary to determine whether a coerced voter complies with the demands [even when the adversary provides the keys or watches the process]" \cite{haines2019surveying}.
Coercion-resistant protocols are implemented by giving voters \emph{deniable} ways to interact with the system.
This can include issuing fake, but indistinguishable credentials or alternative transcripts so that any behaviour observed by the coercer cannot be linked to the voter's true intent.

Most E2E-V protocols cannot achieve full coercion resistance in the strongest sense, especially internet voting schemes.
In practice, many protocols aim for weaker threat assumptions or combine E2E verifiability with procedural safeguards such as supervised polling places to reduce real-world risk \cite{255242}.
Even though full coercion resistance is impractical, it remains as a valuable \emph{design benchmark} that guides how protocols handle voter influence, how much information they expose to observers, and what assumptions they make about the voting environmenent.


\section{Overview of Eperio}

\paragraph{Pre-election.}
Let $b$ be the number of ballots and $s$ the number of candidates, and let $u := s\cdot b$.
Let $U$ be the ordered list of all unique markable regions (UMRs) (all ballot-serial-number/position pairs).
Let $S$ be the list of candidates, and let $M$ be the marks list (initially empty) indexed by $U$.

A set of $n$ trustees run distributed key generation to obtain secret shares $(y_1,\dots,y_n)$ and a shared election secret $\kappa$.
Using $\kappa$ as a seed for a PRNG, the trustees construct a candidate \emph{list} (one candidate per UMR) by repeating the list
across ballots and applying a permutation so that there is a randomized candidate ordering for each ballot.

The trustees form the \emph{print table} $P\in(\cdot)^{u\times 3}$ by joining the three aligned lists
\[
P(:,1):=U,\qquad P(:,2):=M, 
P(:,3):=\text{(candidate list)}.
\]
We use the slicing symbol ":" to denote all rows in a given column.
They then generate the \emph{Eperio table} $E\in(\cdot)^{u\times 3\times x}$ consisting of $x$ independent
row-wise shufflings of $P$.
For each instance $k\in[x]$, a fresh random permutation is applied
to all three columns, producing $E(:,1,k),E(:,2,k),E(:,3,k)$.

Finally, the trustees commit to the first and third columns of each instance.
For each $k\in[x]$ they sample randomness $R(1,k),R(2,k)$ and publish commitments
\[
C(1,k) := \mathsf{Commit}(E(:,1,k),R(1,k)),
C(2,k) := \mathsf{Commit}(E(:,3,k),R(2,k)).
\]

\paragraph{Voting.}
Voters receive a paper ballot with the randomized candidate order from $P$.
After marking the ballot, the candidate list is detached and destroyed.
The remaining strip is scanned and retained by the voter as a receipt. 
After the election, the \emph{marks list} $M$ is completed as marked ($1$), unmarked ($0$), or selected for print-audit ($-1$). 
$M$ is then published.

\paragraph{Post-election tally publication.}
At least $t{+}1$ trustees reconstruct $\kappa$ and regenerate $E$ (now shuffling the completed marks list along with the table).
They publish the marks columns for all instances, where each $E(:,2,k)$ is a shuffled version of the public $M$. 
A tally $\tau$ is then computed from any pair $(E(:,2,k),E(:,3,k))$ and published as the asserted outcome. 

\paragraph{Print-Audits.}
Trustees publish a \emph{linkage list} $L$ giving the row index $a$ such that $E(a,1,k)=U_i$ for each audited UMR index $i$ with $M_i=-1$. 
This lets a print auditor locate the relevant row across each shuffled instance and check consistency against their physical ballot, without revealing complete columns.  

\paragraph{Cut-And-Choose Challenge/Response.}
After $\tau$ is posted, the trustees prove consistency via cut-and-choose.
They regenerate $\kappa$ and $E$, then for each instance $k\in[x]$ obtain a public random challenge bit $z\leftarrow\mathrm{PublicCoin}()$.
They publish the corresponding opening randomness and the challenged committed column:
\[
\text{publish } R(z{+}1,k)\ \text{ and }\ E(:,2z{+}1,k).
\]
If $z=0$ they open the commitment to $E(:,1,k)$, and if $z=1$ they open the commitment to $E(:,3,k)$. 
The marks column $E(:,2,:)$ is already public.

\paragraph{Verification.}
There are three levels of verification.

\emph{Receipt Checks (V1):} a voter checks that their scanned receipt’s marked/unmarked pattern matches the corresponding
entries in the published $M$ for that ballot’s UMRs.

\emph{Print-Audit Checks (V2):} using $L(i,k)$ to locate the row $a$ for each audited UMR $i$ and each instance $k$,
the auditor checks the revealed value is consistent with the physical ballot.
If instance $k$ revealed $E(:,1,k)$, then check $E(a,1,k)=U_i$.
Otherwise, check $E(a,3,k)$ equals the candidate printed at that position.

\emph{Gobal Checks (V3):} for each instance $k$, verifiers 
(i) validate the opening $\mathsf{Reveal}(C(z{+}1,k),R(z{+}1,k),E(:,2z{+}1,k))=1$,
(ii) check isomorphism of the revealed pairs:
if $z=0$, check $\{(E(:,1,k),E(:,2,k))\}$ is a permutation of $\{(U,M)\}$. 
If $z=1$, check $\{(E(:,2,k),E(:,3,k))\}$ is consistent with $\tau$.
